\documentclass[11pt]{amsart}

\usepackage[utf8]{inputenc}
\usepackage[T1]{fontenc}
\usepackage{amsmath,amssymb,amsthm,mathtools}
\usepackage[margin=1in]{geometry}
\usepackage{booktabs}
\usepackage{array}
\usepackage{tabularx}
\usepackage{longtable}
\usepackage{listings}
\usepackage{enumitem}
\usepackage{pdflscape}
\usepackage{hyperref}

% ---------- Theorem-like environments ----------
\theoremstyle{definition}
\newtheorem{definition}{Definition}
\newtheorem{example}{Example}
\theoremstyle{remark}
\newtheorem*{remark}{Remark}

% ---------- Macros ----------
\newcommand{\logic}{L}
\newcommand{\sat}{\models}
\newcommand{\nsat}{\not\models}
\newcommand{\adv}{\rightarrowtail}   % adversarial capability relation
\newcommand{\dual}[1]{(#1)^{d}}
\newcommand{\AP}{\mathit{AP}}
\newcommand{\Edit}{\mathit{Edit}}
\newcommand{\M}{\mathcal{M}}

% Cell formatting for tables
\newcolumntype{L}[1]{>{\raggedright\arraybackslash}p{#1}}

\lstdefinestyle{tamarin}{
  basicstyle=\ttfamily\small,
  frame=single,
  breaklines=true,
  columns=fullflexible,
  keepspaces=true,
  showstringspaces=false
}

\title{A Formal Framework for Threat Models and Attacks Across Arbitrary Logics}
\author{}
\date{}

\begin{document}

\maketitle

%====================================================================
\section{Base Logic}
%====================================================================

In this work, we introduce a formal framework for defining threat models and attacks across arbitrary
logics $L = (M, \Phi, \models)$. By instantiating a threat model as a tuple $T = (m, \varphi, \adv)$, we establish a universal
attack definition as any valid transformation $m \adv m'$ that breaks a nominal safety invariant ($m' \nsat \varphi$).
From here, we introduce the notion of structural duality to transform properties ($\varphi$) and construct dual
threat models $T^{d}$. Lastly, a weighted cut-off ratio $\sigma$ can be introduced to quantify the operational impact of
a successful attack. % NOTE: the source PDF's extraction of this sentence was garbled around the symbol sigma; please check the intended symbol/subscript.
It extends binary invariant violations into a measurable metric that evaluates how
severely an attack degrades valid system execution.

To begin this exploration, the definitions of a logic, threat model, and attack must first be flushed out.
This provides the universal baseline required to analyze security violations across any system.

\begin{definition}
A \emph{logic} is a tuple, $L = (M, \Phi, \models)$, where\ldots
\begin{itemize}
  \item $M$ is a set of models, the things the logic talks about, e.g.\ processes in LTL or worlds in DEL;
  \item $\Phi$ is a set of properties you're allowed to write down about the models;
  \item $\models$ is a satisfaction relation, so that $m \models \varphi$ means $\varphi$ is true about $m$, for any $m \in M$ and $\varphi \in \Phi$.
\end{itemize}
Formally: $\models\ \subseteq M \times \Phi$.
\end{definition}

\begin{definition}
A \emph{threat model} in a logic $L = (M, \Phi, \models)$ is a tuple $(m, \varphi, \adv)$ where\ldots
\begin{itemize}
  \item $m \in M$ is the baseline model, representing the uncompromised system;
  \item $\varphi \in \Phi$ is a safety invariant satisfied by $m$, i.e., $(m \models \varphi)$;
  \item $\adv\ \subseteq M \times M$ is an adversarial capability relation, where $m \adv m'$ denotes that the adversary can
  change model $m$ to model $m'$. We do not assume that $\adv$ is associative, since this would rule out
  notions of bounded capability.
\end{itemize}
\end{definition}

\begin{definition}
Given a logic $L = (M, \Phi, \models)$ and a corresponding threat model $T = (m, \varphi, \adv)$, an \emph{attack} is a modified
model $m' \in M$ where $m \adv m'$, such that $m'$ violates $\varphi$. Formally:
\[
  (m \adv m') \wedge (m' \nsat \varphi).
\]
\end{definition}

\begin{remark}
This logic could also instead be expressed in Hoare Logic where the threat model is a tuple $(p, \varphi, \Edit)$
where\ldots
\begin{itemize}
  \item $p \in P$ is the program, representing the uncompromised source code;
  \item $\varphi \in \Phi$ is a safety invariant satisfied by $p$, i.e., $(p \models \varphi)$;
  \item $\Edit \subseteq P \times P$ is an edit relation, where $p' \in \Edit(p)$ denotes that the adversary can edit
  program $p$ into modified program $p'$.
\end{itemize}
\end{remark}

\begin{remark}
Philosophically, the above Hoare Logic framework shares a great deal of structure with O'Hearne's work,
as both utilize under-approximate reasoning to establish the presence of specification violations. Despite
this, the underlying mechanics of the above and O'Hearne's work differ:
\begin{itemize}
  \item Incorrectness Logic proves error reachability over runtime execution traces ($[P]\ C\ [Q]$), whereas
  the above logic proves specification violations over program transformations ($\Edit \subseteq P \times P$).
  \item Incorrectness Logic targets a single static program but the above framework targets a family of
  modified program variants generated via code transformation ($\{p\}\ \Edit\ \{p'\}$).
  \item Incorrectness Logic defines a strict statement-level proof system with fixed inference rules for
  imperative execution, while the above framework operates at a far more theoretic level, utilizing
  syntax tree mapping to allow for any logical structure.
\end{itemize}
These two frameworks have a lot in common syntactically, but differ fundamentally in function.
\end{remark}

\begin{example}[Defining a Logic for Several Popular Models]
\label{ex:logics}
\end{example}

\begingroup
\small
\renewcommand{\arraystretch}{1.3}
\begin{longtable}{L{2.2cm} L{2.4cm} L{3.3cm} L{2.8cm} L{4.0cm}}
\toprule
\textbf{Logic} & \textbf{Models} $(M)$ & \textbf{Properties} $(\Phi)$ / Formulas built from/bound by: & \textbf{Satisfaction Relation} $(\models)$ & \textbf{Threat Model Tuple} $T = (m, \varphi, \adv)$ \\
\midrule
\endhead
Linear Temporal Logic (LTL) & Set of execution traces & atomic propositions and operators $(G, F, X, U)$ & trace $w$ satisfies temporal formula $\varphi$ &
$m$: execution trace system\newline $\varphi$: LTL temporal safety property\newline $\adv$: trace operator \\
\addlinespace
Computation Tree Logic (CTL) & Set of Kripke structures & path quantifiers $(A, E)$ and linear operators $(X, F, G, U)$ & state $s_0$ in $S$ satisfies formula $\varphi$ &
$m$: Kripke structure\newline $\varphi$: branching-time path property $(AG, EF, AX, \text{etc.})$\newline $\adv$: Transition edge or state-labeling mutation \\
\addlinespace
Dynamic Epistemic Logic (DEL) & Set of pointed epistemic Kripke models & epistemic modalities $(K_a \varphi)$ and action updates & world $w$ in model $\mathfrak{M}$ satisfies $\varphi$ &
$m$: pointed epistemic Kripke model\newline $\varphi$: epistemic or announcement property $([\alpha]\psi,\ K_a \psi)$\newline $\adv$: Action model or announcement update modification \\
\addlinespace
Probabilistic LTL (PCTL) & Set of DTMCs or MDPs & threshold operators $(\mathbb{P}_{\geq \mathbf{p}}(\varphi))$ & probability of satisfying path property $\varphi$ meets bound $\mathbf{p}$ &
$m$: Markov Chain\newline $\varphi$: probabilistic bound formula $(\mathbb{P}_{\bowtie \mathbf{p}}(\psi))$\newline $\adv$: Probability weight perturbation operator \\
\addlinespace
Signal Temporal Logic (STL) & Set of continuous time signals & dense-time temporal operators $(G_{[a,b]})$ & signal trajectory $s$ satisfies metric property $\varphi$ at time $t$ &
$m$: continuous-time signal trajectory\newline $\varphi$: temporal logic property $(G_{[a,b]},\ F_{[a,b]})$\newline $\adv$: Continuous signal amplitude \\
\addlinespace
Turing Machine & set of execution traces & Safety properties $(\varsigma)$ & trace $w$ satisfies safety property $\varphi$ &
$m$: execution trace system\newline $\varphi$: safety property\newline $\adv$: transition rule mutation \\
\addlinespace
Micro-architectural Side-Channel Attacks on CPU & set of microarchitectural transition systems & security properties specifying non-interference or timing indistinguishability & processor model $m$ securely isolates secrets $\varphi$ &
$m$: a secure CPU model\newline $\varphi$: security property\newline $\adv$: adversarial microarchitectural change \\
\bottomrule
\end{longtable}
\endgroup

%====================================================================
\section{Duality}
%====================================================================

A common theme in many logics is that some properties are stronger than others. This can be formalized
in many ways, for example, in a logic over program traces, a stronger property is one that holds in a
smaller number of traces (it is stronger in the sense that it is more precise). Often, a logic will be
expressed using variables, models, and quantifiers, of some kind. In such cases, it is common for one
quantifier to be the dual of another. More precisely, if we have quantifiers $J$ and $K$, we say they are duals
if $J = \neg K \neg$, (e.g.\ $\forall = \neg \exists \neg$). In this next section, we discuss how this dual relationship
matters for attacks---it often captures ideas around the strength of an attack, such as if it always works or
sometimes works, and in different logical languages can take different forms, as captured in Example~2.1.
We generalize the concept, mathematically, over arbitrary logics, and then we use our generalization to
instantiate what it means for several interesting logics to show its usefulness.

\begin{remark}
By the definition of a dual, one might assume that a dual is either particular to a logical language or that
multiple sets of duals cannot exist within the same logical language. Neither is true, for example in First
Order Logic (FOL), you have both the duals ``and'' and ``or'' and also ``for every'' and ``there exists''; both
sets exist within the same framework, and both sets are common to many logical languages.
\end{remark}

\begin{definition}
Let $\Phi$ be the set of full formulas generated inductively over atomic variables $\AP$ and operators $\Omega$. The
\emph{structural (De Morgan) dual} $(\cdot)^{d} : \Phi \to \Phi$ is defined over the primary constructor encountered at the root
of a formula's abstract syntax tree, satisfying $\big((\varphi)^{d}\big)^{d} = \varphi$:
\begin{itemize}
  \item Literals: $(\mathsf{p})^{d} = \neg \mathsf{p}$ and $(\neg \mathsf{p})^{d} = \mathsf{p}$
  \item Compound Formulas: $\big(\omega(\psi_1, \ldots, \psi_n)\big)^{d} = \neg\big(\omega(\neg\psi_1, \ldots, \neg\psi_n)\big)$
\end{itemize}
\end{definition}

\begin{definition}
A \emph{dual threat model} replaces the target property $\varphi$ with its structurally transformed dual formula $(\varphi)^{d}$,
preserving the baseline model $m$ and capability relation $\adv$:
\[
  T^{d} = \big(m, (\varphi)^{d}, \adv\big).
\]
\end{definition}

\begin{definition}
Given a dual threat model $T^{d} = (m, (\varphi)^{d}, \adv)$ and the candidate space $C(m) = \{m' \in M \mid m \adv m'\}$, a
\emph{realized dual attack} occurs when a transitioned model $m' \in C(m)$ satisfies the dual formula:
\[
  m' \models (\varphi)^{d}.
\]
\end{definition}

\begin{example}[Languages and Their Duals]
\label{ex:duals}
\end{example}

\begingroup
\small
\renewcommand{\arraystretch}{1.4}
\begin{longtable}{L{1.8cm} L{2.6cm} L{3.4cm} L{3.4cm} L{3.4cm}}
\toprule
\textbf{Logic} & \textbf{Duals} & \textbf{Safety Specification} $(\varphi)$ Examples & \textbf{Dual Adversarial Goal} $(\varphi)^{d}$ Examples & \textbf{Impact Examples} \\
\midrule
\endhead
LTL &
$(G, F)$\newline $(\wedge, \vee)$\newline $(\Rightarrow, \wedge)$ &
$G(\mathit{Valid} \Rightarrow F\,\mathit{Secure})$ &
$F(\mathit{Valid} \wedge G\neg \mathit{Secure})$ &
Transforms a safety invariant into a trace where safety degrades. \\
\addlinespace
CTL &
$(A, E)$\newline $(G, F)$ &
$AG(\mathit{State} = \mathit{Normal})$ &
$AG(\mathit{State} \neq \mathit{Normal})$ &
demonstrates the existence of at least one reachable compromised branch \\
\addlinespace
DEL &
$([\alpha], \langle\alpha\rangle)$\newline $(K_a, \hat{K}_a)$\newline $(\wedge, \vee)$ &
$[\mathit{Read}]\,\mathit{SecretSafe}$ &
$\langle \mathit{Read} \rangle \neg \mathit{SecretSafe}$ &
proves an execution path exists where the adversary successfully gains unauthorized information. \\
\addlinespace
PCTL &
$(G, F)$\newline $(\geq, <)$\newline $(>, \leq)$ &
$\mathbb{P}_{\geq 0.99}(G\,\mathit{StateOK})$ &
$\mathbb{P}_{<0.99}(F\,\mathit{StateFail})$ &
Converts a reliable bound into a likelihood that an adversary can trigger a failure. \\
\addlinespace
STL &
$(G, F)$\newline $(\geq, <)$\newline $(>, \leq)$ &
$G_{[0,\mathsf{T}]}(x(t) \geq \mathit{Threshold})$ &
$F_{[0,\mathsf{T}]}(x(t) < \mathit{Threshold})$ &
Injects an adversarial change breaking continuous signal safety at a specific moment in time. \\
\bottomrule
\end{longtable}
\endgroup

%====================================================================
\section{Tamarin}
%====================================================================

To enable the use of a wide variety of formal verification tools, one can actually transfer the logics or
modalities that these tools work under directly into Hoare-style triples. Rather than viewing verification
tools as isolated external engines, our framework lifts their state-transition mechanics directly into our
deductive Hoare logic structure. For example,

\begin{example}[LTL to Tamarin to Hoare Logic]
\label{ex:tamarin}
Within Tamarin we have pretty clear equivalencies on how the Tamarin Framework can become
Hoare-logic triples.
\begin{itemize}
  \item Preconditions and States ($P$): Program variables are represented as multisets of active facts
  \item adversarial capability relations ($\adv$): An adversarial program mutation is mapped to a rewrite rule
  consuming pre-state facts and producing post-state facts.
  \item Invariant violations ($p' \nsat \varphi$): Postcondition violations are evaluated as trace properties (lemma)
\end{itemize}

Using an example LTL formula:
\[
  \{\mathit{balance} \geq 0\}\ \mathtt{if}\ (x \leq \mathit{balance})\ \mathtt{then}\ \mathit{balance} := \mathit{balance} - x\ \mathtt{else}\ \mathtt{skip}\ \{\mathit{balance} \geq 0\}
\]
we can generate the Tamarin file:

\begin{lstlisting}[style=tamarin]
theory AccountMutation
begin

builtins: multiset
functions: minus/2

// 1. Initial State Rule
rule Init_Account:
  [ In(~initial_balance) ]
  --[ Setup(~initial_balance) ]->
  [ Balance(~initial_balance) ]

// 2. Adversarial Program Mutation Rule (Edit)
rule Malicious_Withdraw_Edit:
  [ Balance(bal), In(x) ]
  --[ StateTransition(bal, x) ]->
  [ Balance(minus(bal, x)) ]

// 3. Safety Invariant Trace Lemma
// Note: Tamarin parses formulas using standard logical connectives (=, not, &, |, ==>)
lemma safety_invariant_non_negative:
  " All bal x #i. StateTransition(bal, x) @ i ==> not (bal = x) "

end
\end{lstlisting}

Running this through Tamarin results in our valid attack, successfully breaking our security property:

\begin{lstlisting}[style=tamarin]
theory AccountMutation begin

Message theory
Multiset rewriting rules (4)
Tactic(s)
Raw sources (5 cases, deconstructions complete)
Refined sources (5 cases, deconstructions complete)

add lemma
lemma safety_invariant_non_negative:
  all-traces
  "forall bal x #i.
    (StateTransition( bal, x ) @ #i) ==> (not(bal = x))"
edit lemma or delete lemma
by sorry
add lemma

end
\end{lstlisting}

This corresponds to our Edit relation:
\[
  \{\mathit{balance} \geq 0\}\ \Edit\ \{\mathit{balance} < 0\}
\]
Which is inherently what we got earlier from our LTL formula when making the transfer directly into our
logical system.
\end{example}

\begin{example}[Model Checkers and Their Use Within the Overall Framework]
\label{ex:checkers}
Each row is derived from the tool's own documentation or paper (cited in the first column), not from
the framework. The last column states how the tool instantiates $(M, \Phi, \models)$ and the adversarial
capability relation $\adv$. Where a tool has no adversary of its own, the capability relation must be supplied
by the user as part of the model.
\end{example}

\begin{landscape}
\begingroup
\footnotesize
\setlength{\tabcolsep}{3pt}
\renewcommand{\arraystretch}{1.25}
\begin{longtable}{L{1.7cm} L{4.0cm} L{3.5cm} L{4.2cm} L{4.4cm} L{3.6cm}}
\toprule
\textbf{Tool} & \textbf{Models} $(M)$ & \textbf{Properties} $(\Phi)$ & \textbf{Adversary} & \textbf{Analysis} & \textbf{Instantiation of $T = (m, \varphi, \adv)$} \\
\midrule
\endfirsthead
\toprule
\textbf{Tool} & \textbf{Models} $(M)$ & \textbf{Properties} $(\Phi)$ & \textbf{Adversary} & \textbf{Analysis} & \textbf{Instantiation of $T = (m, \varphi, \adv)$} \\
\midrule
\endhead
Tamarin \cite{tamarin} &
Multiset-rewriting rules defining a labeled transition system over protocol state, messages, and
the adversary's symbolic knowledge. &
Trace properties (first-order formulas over timestamped action facts) and observational
equivalence. &
Also rules in the same language; standard instantiation is Dolev--Yao deduction. &
Constraint solving: exhaustive symbolic search for a trace violating the property. &
$m$: rule set; $\varphi$: lemma; $\adv$: edit a rule; attack: an edited rule set's trace violating the lemma. \\
\addlinespace
SPIN \cite{spin} &
Promela processes and channels (rendezvous or buffered); execution interleaves enabled statements. &
Assertions, invalid end states, progress/acceptance labels, and never claims (LTL). &
None built in; nondeterministic scheduling explores all interleavings, with the environment
modeled as ordinary processes. &
Exhaustive DFS with partial-order reduction; bitstate hashing for large models (sound against
false positives, may miss errors). &
$m$: process composition; $\adv$: replacing or adding a process, as in the attacker synthesis of \cite{vonhippel}. \\
\addlinespace
ProVerif \cite{proverif} &
Protocols in applied $\pi$-calculus, with cryptographic primitives as constructor/destructor
rewrite rules. &
Secrecy (reachability of the attacker-knowledge predicate), correspondence assertions, and
diff-equivalence. &
Symbolic Dolev--Yao attacker with perfect cryptography, encoded as Horn clauses over the predicate
$\mathit{att}(M)$ rather than a separate component. &
Horn-clause resolution, abstracting protocol-level constraints to handle unboundedly many
sessions; may not terminate and is incomplete (false attacks). &
$m$: process; $\varphi$: query; $\adv$: fixed by the tool (Dolev--Yao); attack: a derivable fact that violates the query. \\
\addlinespace
nuXmv \cite{nuxmv} &
Finite- and infinite-state synchronous transition systems, with Integer, Real, and
uninterpreted-function types for the infinite case. &
Invariants, and LTL and CTL properties inherited from NuSMV. &
None built in. &
Finite-state: SAT-based algorithms, interpolation, IC3, $k$-liveness. Infinite-state: SMT-based
BMC, $k$-induction, interpolation, IC3, and CEGAR abstraction. &
$m$: transition system; $\adv$: mutation of the transition relation or variable assignments, supplied by the user. \\
\addlinespace
PRISM \cite{prism} &
DTMCs, CTMCs, MDPs, and PTAs, each optionally augmented with rewards. &
PCTL (DTMC, MDP), CSL (CTMC), LTL, PCTL*, quantitative queries ($P{=}?$), and reward operators. &
None built in; nondeterminism in MDPs is resolved by minimizing or maximizing probabilities. &
Matrix--vector multiplication, linear equation solving, Prob0/Prob1 precomputation, iterative
solvers, value iteration, and BDD/MTBDD symbolic storage. &
$m$: DTMC; $\adv$: perturbing the transition matrix, as in \cite{oakley}, which uses PRISM as the inner probability engine. \\
\addlinespace
HyTech \cite{hytech} &
Hybrid automata composed as a product: nodes carry invariants and differential inclusions,
transitions carry guards, resets, and synchronization labels; linear dynamics (stopwatches, skewed
clocks, bounded rates). &
State-set reachability queries; the logic is not named in the paper. &
None built in. &
Symbolic model checking over polyhedra (termination not guaranteed), plus parametric analysis and
diagnostic error traces. &
$m$: hybrid automaton; $\adv$: edits to rates, guards, and resets, supplied by the user. \\
\addlinespace
Alloy \cite{alloy} &
Relational structures over a finite, user-chosen scope: fixed signatures and the relations
(``fields'') declared over them. &
Facts (constraints every valid instance must satisfy) and assertions (claimed consequences),
both relational first-order formulas. &
None built in; the natural $\adv$ is the universal relation over every instance satisfying the
same facts, since \texttt{check} asks whether \emph{every} such instance also satisfies the assertion. &
Translation to boolean logic, solved by SAT, bounded by the declared scope (the ``small scope
hypothesis''): no counterexample within scope is not a proof none exists at a larger one. &
$m$: a relational instance; $\varphi$: an assertion; $\adv$: any other instance satisfying the same facts;
attack: a counterexample instance violating the assertion. \\
\addlinespace
UPPAAL \cite{uppaal} &
Networks of timed automata: locations with clock invariants, edges with guards and resets,
synchronized over channels; every clock advances at a fixed rate. &
A decidable fragment of TCTL: reachability and safety queries over the automata network. &
None built in; the natural $\adv$ edits a guard or invariant threshold (a timing margin), the
only lever available to compromise a timed system's safety. &
Symbolic reachability over zones (regions of clock-value space), explored on the fly. &
$m$: a network configuration (its guards/invariants); $\varphi$: a safety query (a location is never
reached); $\adv$: a threshold edit; attack: an edit after which the forbidden location becomes reachable. \\
\addlinespace
Margrave \cite{margrave} &
A policy as a mapping from requests to decisions: an ordered list of rules combined ``first
applicable.'' &
First-order queries over per-rule $R_{\mathit{matches}}$/$R_{\mathit{applies}}$ formulas, e.g.\
whether a given request is ever permitted. &
None built in; the paper's own ``change impact'' analysis -- over edits to the rule list -- is
exactly the $\adv$ used here: insert, delete, flip, or reorder a rule. &
Compilation to first-order relational formulas, solved by the Kodkod/Alloy SAT backend over a
bounded universe. &
$m$: an ordered rule list; $\varphi$: a fixed request must decide ``deny''; $\adv$: a single rule edit;
attack: an edit after which the request decides ``permit.'' \\
\bottomrule
\end{longtable}
\endgroup
\end{landscape}

%====================================================================
\section{Measure Theory Extension}
%====================================================================

\begin{definition}
A \emph{probability measure space} is a triple $(\M, \Sigma, \mu)$ where:
\begin{itemize}
  \item $\M$ is a non-empty set of models;
  \item $\Sigma \subseteq \mathcal{P}(\M)$ is a $\sigma$-algebra over $\M$ containing $\M$ and closed under complementation and countable
  unions;
  \item $\mu : \Sigma \to [0, 1]$ is a probability measure satisfying $\mu(\emptyset) = 0$, $\mu(\M) = 1$, and countable additivity
  for any sequence of pairwise disjoint sets $(M_i)_{i=0}^{\infty} \subset \Sigma$.
\end{itemize}
\end{definition}

\begin{definition}
Let $\M$ be the set of models for the logic, equipped with a probability measure space $(\M, \Sigma, \mu)$
normalized such that $\mu(\M) = 1$. Each subset of models $M \in \Sigma$ is assigned a probability mass $\mu(M)$
representing its relative execution share, frequency, or occurrence weight.
\end{definition}

\begin{definition}
The \emph{severity} of an attack is quantified by the weighted model cut-off ratio $\sigma \in [0, 1]$, defined as the
proportional loss of measure between the model set $M$ and the attacked model set $M'$:
\[
  \sigma = 1 - \frac{\mu(M')}{\mu(M)}.
\]
\begin{itemize}
  \item $\sigma = 0$: Nominal operation
  \item $0 < \sigma < 1$: Some Branches are severed, but the program remains executable
  \item $\sigma = 1$: System is fully compromised
\end{itemize}
\end{definition}

%====================================================================
\section{Comparisons}
%====================================================================

\begin{example}[Comparison to Other Attack Synthesis Frameworks]
\label{ex:comparison}
Each column is derived from the cited paper's own definitions (section and definition numbers are given
where the paper provides them). Where a paper does not define a component, the cell says so. Entries
marked ``(our encoding)'' are not stated in the paper but follow from its definitions.
\end{example}

\begin{landscape}
\begingroup
\footnotesize
\setlength{\tabcolsep}{3pt}
\renewcommand{\arraystretch}{1.25}
\begin{longtable}{L{3.0cm} L{4.6cm} L{4.6cm} L{4.6cm} L{4.6cm}}
\toprule
\textbf{Component} & \textbf{Oakley et al.\ \cite{oakley}} & \textbf{von Hippel et al.\ \cite{vonhippel}} & \textbf{Bailluet et al.\ \cite{bailluet}} & \textbf{Wardega et al.\ \cite{wardega}} \\
\midrule
\endfirsthead
\toprule
\textbf{Component} & \textbf{Oakley et al.\ \cite{oakley}} & \textbf{von Hippel et al.\ \cite{vonhippel}} & \textbf{Bailluet et al.\ \cite{bailluet}} & \textbf{Wardega et al.\ \cite{wardega}} \\
\midrule
\endhead
Logic $L = (M, \Phi, \models)$ &
PCTL* over DTMCs. MDPs are covered under a deterministic memoryless policy, which induces a DTMC. &
LTL over processes. $\mathsf{P} \models \varphi$ means every computation of $\mathsf{P}$ satisfies $\varphi$. &
Quantifier-free formulas over bit-vectors and arrays (QF\_ABV) describing machine states and gadget transitions. &
No logic is defined; satisfaction is given by the paper's definitions of plans, deviations, and detection (Defs.~1--6). The paper does not use an epistemic logic. \\
\addlinespace
Model $m \in M$ &
DTMC $C = (S, s_0, \mathbf{P}, \mathit{AP}, L)$. &
Composition $\mathsf{P} \parallel Q_0 \parallel \dots \parallel Q_n$ of Kripke-structure processes with inputs and outputs, composed asynchronously with rendezvous. &
Initial machine state $s_0 = \langle \mathbf{p}, \mathbf{u} \rangle$ (registers and byte memory), split into attacker-predictable $\mathbf{p}$ and unpredictable $\mathbf{u}$, together with a library of gadgets. &
A graph $G = (V, E)$ of locations, robots $R$, a MAPF plan $x$, and the prefix $\alpha(t)$ of $x$ announced by the central entity at time $t$. \\
\addlinespace
Safety invariant $\varphi \in \Phi$ &
$\mathbb{P}_{\geq \mathbf{p} - \delta}(\psi)$, where $\psi$ is a PCTL* path formula and $\mathbf{p} = \Pr_C(\psi)$ (robustness, Def.~6). &
An LTL property $\varphi$ with $\mathsf{P} \parallel Q_0 \parallel \dots \parallel Q_n \models \varphi$ before the attack. &
Not an invariant in the paper; (our encoding) $\varphi_n = \neg\, GS_n$: no chain of $n$ gadgets robustly reaches the goal specification $S$ from assumptions $A$. &
Horizon-limiting announcements (Def.~6): for every forbidden deviation, some continuation of $\alpha(t)$ makes the reports differ, so the attacker cannot be certain of going undetected. \\
\addlinespace
Adversarial capability relation $\adv\ \subseteq M \times M$ &
$C \adv C'$: same $S$, $s_0$, $L$, and a matrix $\mathbf{P}'$ with $\|\mathbf{P} - \mathbf{P}'\|_{\max} \leq \varepsilon$, changing only selected states (SS) or transitions (ST); structure-preserving variants keep zero entries zero. &
$\mathsf{P} \parallel Q \adv \mathsf{P} \parallel A$: each $Q_i$ replaced by a deterministic $A_i$ with the same interface and no atomic propositions; the recovery variant uses a finite DAG ending in $Q_i$. &
Rewrites the code pointer and controlled memory data, moving between states only via gadget transitions $g(s, s') = g_{\mathit{pre}}(s) \wedge g_{\mathit{post}}(s, s')$ (NX on, CFI off). &
A compromised robot $i$ changes only its own positions on $(s, f)$ and falsifies its self-reports (Def.~3), knowing $\alpha(t)$ but not which other robots are compromised. \\
\addlinespace
Attack $(m \adv m') \wedge (m' \nsat \varphi)$ &
A perturbation $C'$ with $\Pr_{C'}(\psi) < \Pr_C(\psi) - \delta$. &
An attacker $A$ with $\mathsf{P} \parallel A \nsat \varphi$ (an $\exists$-attacker if it succeeds on some runs). &
An assignment to $\mathbf{p}$ and a chain of $n$ gadgets for which $S$ holds, for all values of $\mathbf{u}$ (Def.~1, robust $n$-gadget satisfiability). &
A deviation that enters $V_{\mathit{forbidden}}$ (Def.~4) while the central entity receives exactly the expected reports. \\
\addlinespace
Attacker goal $m' \models (\varphi)^d$ &
Worst-case attack: the $C'$ maximizing $\Pr_C(\psi) - \Pr_{C'}(\psi)$, i.e.\ $C' \models \mathbb{P}_{< \mathbf{p} - \delta}(\psi)$. &
A $\forall$-attacker: $\mathsf{P} \parallel A \models \neg\varphi$, so every behavior violates $\varphi$. &
The specification $S$ itself: $S(s_n)$ holds after $n$ gadgets. &
Stealthy attacker: a forbidden undetected deviation exists for every continuation of $\alpha(t)$, which is the negation of the horizon-limiting property. A bold attacker tries without that guarantee. \\
\addlinespace
Verifier &
Optimization with SLSQP over PRISM, or with PARAM for a closed-form parametric solution. &
Reduction to model checking: check $\mathsf{P} \parallel \mathrm{DAISY}(Q_0) \parallel \dots$ against $\varphi$ with SPIN, then project the counterexample onto each $Q_i$ (tool: KORG). Sound and complete for $\exists$ attackers. &
One SMT query over bit-vectors and arrays with the Boolector solver, using a taint encoding to eliminate $\forall \mathbf{u}$; chains are searched by increasing length, with gadget sampling for scale. &
Algorithm~1, a non-deterministic co-observation enumeration: sound but incomplete, with cost $O(|R||V|)$, checked per robot. \\
\addlinespace
Severity / impact (cut-off ratio $\sigma$) &
Not defined as a ratio; the paper maximizes the probability drop $\Pr_C(\psi) - \Pr_{C'}(\psi)$. &
Not defined; attackers are classified as $\exists$ versus $\forall$, with or without recovery. &
Not defined; the chain either exists or not, and solving time is reported. &
Not defined as a ratio; the paper reports the share of scenarios verified secure, bold-attack success and miss rates, and lookahead. \\
\bottomrule
\end{longtable}
\endgroup
\end{landscape}

%====================================================================
\section{Mechanization}
%====================================================================

Sections~1--4 define the framework abstractly: a logic, a threat model, an attack, a dual threat
model, a realized dual attack, and a severity ratio. Realizing any one of these for a particular
logic $L$ and threat model $T$ requires only two computable ingredients: the satisfaction relation
$\models$, and the capability relation $\adv$. We provide an open-source reference implementation,
\texttt{attack\_synth.py} \cite{attacksynth}, that takes exactly these two ingredients and performs
an automated search for a witness of Definition~3 (an attack) or Definition~6 (a realized dual
attack), computing Definition~9's severity automatically whenever a measure is supplied. This
section describes that implementation and what mechanizing the framework does and does not buy us.

\begin{definition}
\leavevmode\\
Given a logic $L = (M, \Phi, \models)$ and a threat model $T = (m, \varphi, \adv)$ with computable
$\models$ and $\adv$, the \emph{realized-attack search problem} asks for some $m' \in C(m) =
\{m' \in M \mid m \adv m'\}$ with $m' \nsat \varphi$ (or, dually, $m' \models (\varphi)^{d}$). This
problem is recursively enumerable whenever $C(m)$ is computably enumerable, and decidable whenever
$C(m)$ is additionally finite.
\end{definition}

\begin{remark}
Because $\adv$ is not assumed associative (Definition~2), a path $m \adv m_1 \adv \cdots \adv m'$ of
$k$ composed edits is a choice the user makes about how much capability to grant the adversary, not
a property derivable from $\adv$ alone. Our implementation exposes this choice directly as a
parameter \texttt{max\_steps}, defaulting to $k = 1$ (a single edit, exactly Definition~3).
\end{remark}

The implementation is organized around three objects, corresponding directly to the tuples above:

\begin{itemize}
  \item \texttt{Logic(name, verifier)} -- $\models$, with $M$ and $\Phi$ implicit, encoded as the
  boolean function $f(m, \varphi) = \mathsf{True}$ iff $m \models \varphi$.
  \item \texttt{ModelType(universe, neighbors, key, perturb)} -- the candidate space the adversary
  can reach. $\adv$ is likewise encoded as a boolean function
  \texttt{edit\_relation}$(m, m') = \mathsf{True}$ iff $m \adv m'$, applied either by filtering an
  over-approximate \texttt{universe} (exact, for small enumerable $M$) or by generating
  already-restricted \texttt{neighbors} (for $M$ too large or continuous to enumerate, e.g.\ an
  $\varepsilon$-ball of perturbed DTMCs or a gadget library).
  \item \texttt{ThreatModel(logic, m, phi, model\_type, edit\_relation, measure, max\_steps)} -- the
  tuple $T = (m, \varphi, \adv)$ itself. Constructing one checks $m \models \varphi$ (Definition~2's
  precondition) and raises an error if it fails, rather than searching from an invalid baseline.
\end{itemize}

Given a \texttt{ThreatModel}, \texttt{synthesize()} searches $C(m)$ for a witness: exhaustive
breadth-first up to \texttt{max\_steps} by default, or, when a measure $\mu$ is supplied, best-first
by severity ($\sigma$, Definition~9) so a budget still finds a high-severity attack quickly on a
space too large to exhaust. A further \texttt{robust} flag asks for an attack holding for every
\texttt{perturb}$(m')$ of a candidate -- the ``for all unpredictable inputs'' quantifier of Bailluet
et al.\ \cite{bailluet}'s robust gadget chains (Example~\ref{ex:comparison}).

\begin{table}[h]
\centering
\small
\renewcommand{\arraystretch}{1.3}
\begin{tabular}{L{4.2cm} L{4.2cm}}
\toprule
\textbf{Definition} & \textbf{Code construct (\texttt{attack\_synth.py})} \\
\midrule
Logic $L = (M, \Phi, \models)$ & \texttt{Logic(name, verifier)} \\
Threat model $T = (m, \varphi, \adv)$ & \texttt{ThreatModel(...)}; raises unless $m \models \varphi$ \\
Attack $(m \adv m') \wedge (m' \nsat \varphi)$ & \texttt{ThreatModel.synthesize(goal="violate")} \\
Structural dual $(\cdot)^{d}$ & \texttt{structural\_dual(phi)} \\
Dual threat model $T^{d}$ & \texttt{DualThreatModel(base, phi\_d)} \\
Realized dual attack $m' \models (\varphi)^{d}$ & \texttt{DualThreatModel.realized()} \\
Probability measure space $(\M, \Sigma, \mu)$ & \texttt{MeasureSpace(mu)} \\
Severity $\sigma = 1 - \mu(M')/\mu(M)$ & \texttt{severity(mu, m, m')} \\
\bottomrule
\end{tabular}
\caption{Each Definition in Sections~1--4 and the code construct realizing it literally.}
\label{tab:conformance}
\end{table}

\begin{example}[Instantiated Model Classes]
\label{ex:instances}
\leavevmode\\
Without modifying the search engine, the implementation has been instantiated against all eight
logics of Example~\ref{ex:logics}, all four frameworks of Example~\ref{ex:comparison}, and all nine
tools of Example~\ref{ex:checkers} -- twenty-one instances in total, each reusing that logic's,
paper's, or tool's own model class and adversary, with no change to \texttt{ModelType} or the search
procedure: only new \texttt{verifier}/\texttt{edit\_relation} pairs.
\end{example}

\begin{remark}
Mechanization is checked, not asserted: an accompanying test suite reconstructs each Definition's
precondition and conclusion as an executable assertion -- that $(\cdot)^{d}$ is an involution, that
a found attack genuinely satisfies $(m \adv m') \wedge (m' \nsat \varphi)$ against the raw
$\models$/$\adv$ (not merely trusted from the search), that a realized dual attack genuinely
satisfies $m' \models (\varphi)^{d}$, that $\sigma$ is clamped to $[0,1]$ and attains all three cases
of Definition~9, and that the \texttt{ThreatModel} API and the lower-level search engine agree
exactly on a shared instance. All tests pass, including against a model class (Margrave's
first-applicable rule ordering) added after this API was written.
\end{remark}

\paragraph{Implications.} Mechanizing the framework does not add content to it -- every result above
follows from Sections~1--4 alone -- but it changes what ``this framework applies to arbitrary
logics'' means from a claim to a falsifiable one: Example~\ref{ex:instances} is twenty-one chances
for the claim to have failed, and it did not. The generality is bounded exactly where
Definitions~1--2 say it should be: the engine supplies no $\models$ or $\adv$ of its own, so it is a
search substrate, not an oracle, and a ``no attack found'' result is only as strong as the fidelity
of the supplied model and the completeness of the search budget -- the same caveat already attached
to several tools in Example~\ref{ex:checkers} (ProVerif's incompleteness, HyTech's non-terminating
symbolic reachability). What mechanization adds operationally is severity as a common currency:
none of the four frameworks in Example~\ref{ex:comparison} define a severity ratio, reporting
instead a success/failure flag, a probability drop, or a solving time specific to their own logic.
Definition~9's $\sigma \in [0,1]$ lets attacks from entirely different logics -- a DTMC perturbation,
a firewall misconfiguration, a timed-automaton race -- be ranked on one scale, which no individual
logic or tool provides on its own.

\begin{thebibliography}{9}
\bibitem{oakley} L.\ Oakley, A.\ Oprea, S.\ Tripakis. \emph{Adversarial Robustness Verification and Attack Synthesis in Stochastic Systems}. arXiv:2110.02125.
\bibitem{vonhippel} M.\ von Hippel, C.\ Vick, S.\ Tripakis, C.\ Nita-Rotaru. \emph{Automated Attacker Synthesis for Distributed Protocols}. arXiv:2004.01220.
\bibitem{bailluet} N.\ Bailluet, E.\ Fleury, I.\ Puaut, E.\ Rohou. \emph{Nothing is Unreachable: Automated Synthesis of Robust Code-Reuse Gadget Chains for Arbitrary Exploitation Primitives}. USENIX Security Symposium, 2025.
\bibitem{wardega} K.\ Wardega et al. \emph{HoLA Robots: Mitigating Plan-Deviation Attacks in Multi-Robot Systems with Co-Observations and Horizon-Limiting Announcements}. AAMAS 2023; arXiv:2301.10704.
\bibitem{tamarin} The Tamarin Prover Manual. \url{https://tamarin-prover.com/manual/master/book/001_introduction.html}.
\bibitem{spin} The Spin model checker: Promela quick reference. \url{https://web.mit.edu/spin_v6.4.7/Doc/Quick.html}.
\bibitem{proverif} B.\ Blanchet. \emph{The Security Protocol Verifier ProVerif and its Horn Clause Resolution Algorithm}. arXiv:2211.12227.
\bibitem{nuxmv} The nuXmv model checker. \url{https://nuxmv.fbk.eu/} (features page), and the CAV 2014 tool paper.
\bibitem{prism} The PRISM model checker: model types and properties. \url{https://prismmodelchecker.org/doc/whatsinprism.php}.
\bibitem{hytech} T.\ A.\ Henzinger, P.-H.\ Ho, H.\ Wong-Toi. \emph{HyTech: A Model Checker for Hybrid Systems}. CAV 1997.
\bibitem{alloy} \emph{Alloy (specification language)}. \url{https://en.wikipedia.org/wiki/Alloy_(specification_language)}.
\bibitem{uppaal} \emph{UPPAAL Model Checker}. \url{https://en.wikipedia.org/wiki/Uppaal_Model_Checker}.
\bibitem{margrave} T.\ Nelson, C.\ Barratt, D.\ J.\ Dougherty, K.\ Fisler, S.\ Krishnamurthi. \emph{The Margrave Tool for Firewall Analysis}. USENIX LISA, 2010.
\bibitem{attacksynth} \emph{attack-synth}: reference implementation of this paper's framework. \url{https://github.com/samvonhippel/attack-synth}.
\end{thebibliography}

\end{document}
